\documentclass[11pt,a4paper]{article}

\usepackage[T1]{fontenc}
\usepackage[margin=2.5cm]{geometry}
\usepackage{amsmath}
\usepackage{array}
\usepackage{hyperref}

\title{Implementation Plan for the LM Log-Likelihood Backend}
\date{\today}

\begin{document}

\maketitle

\section{Purpose}

This document gives a practical, step-by-step plan for connecting the existing
Languasco--Migliardi (LM) C implementation to the LM-IDNet Python application.
The completed backend will evaluate the parameter-dependent
Dirichlet--Multinomial log-likelihood kernel during model fitting, calibration,
and anomaly scoring.

The work is an integration task. It will not redesign the LM approximation,
the Euler--Maclaurin expansion, the supplied coefficients, or Minka's
fixed-point estimator. The aim is to expose the existing numerical method
through a small, safe native interface and make it satisfy the Python backend
contract already used by the SciPy implementation.

\section{Definition of Done}

The work is complete when all of the following statements are true:

\begin{itemize}

\item selecting \texttt{lm} creates a working backend instead of raising the
placeholder error;

\item the LM backend accepts the same count matrices and Dirichlet parameter
vectors as the SciPy backend;

\item it returns the same row-wise likelihood kernel, within the requested
numerical accuracy;

\item zero-count categories and completely silent windows are handled
correctly;

\item coefficient files are found independently of the current working
directory;

\item repeated calls do not leak native working memory or reuse stale data;

\item native failures become Python application exceptions instead of
terminating the process;

\item the requested precision is passed from application configuration to the
C implementation for training, diagnostics, calibration, and scoring;

\item the targeted numerical tests and the complete test suite pass; and

\item at least one complete \texttt{train}, \texttt{calibrate}, and
\texttt{score} run succeeds with the LM backend.

\end{itemize}

\section{Current Starting Point}

The current codebase already provides most of the surrounding structure:

\begin{itemize}

\item \path{src/lm_idnet/algorithms/log_likelihood.py} defines the common
backend interface and a working SciPy reference implementation.

\item \texttt{LmLogLikelihood} is registered under the configuration value
\texttt{lm}, but it is currently a placeholder.

\item \path{src/lm_idnet/algorithms/lm_native.py} exists and is ready to hold
the Python-to-C adapter.

\item The runtime copy of the C implementation, its headers, coefficient
tables, licence, and attribution are already stored in
\path{src/lm_idnet/algorithms/native/}.

\item The source capsule under \path{lm-algorithm/} contains the original
experimental program and remains useful as reference material. It will not be
used directly by the application and will not be modified during this task.

\item \path{scripts/docker-entrypoint.sh} already compiles
\path{LM_time_lib.c} into \path{LM_time_lib.so} before running a command in
Docker.

\item The validated application configuration already contains
\texttt{precision\_digits} and the selectable likelihood backend.

\item Training stores the chosen backend in the model artifact. Calibration
and scoring reload that backend name from the model.

\item The multinomial coefficient used for complete anomaly scores is already
calculated separately in
\path{src/lm_idnet/algorithms/dirichlet_multinomial.py}.

\end{itemize}

The task should therefore reuse these pieces rather than create a second
estimator, a second likelihood abstraction, or a Python translation of the C
algorithm.

\section{Required Likelihood Semantics}

For one count vector

\[
\mathbf{x}=(x_1,\ldots,x_K),
\qquad
N=\sum_{k=1}^{K}x_k,
\]

and positive Dirichlet parameters

\[
\boldsymbol{\alpha}=(\alpha_1,\ldots,\alpha_K),
\qquad
\alpha_0=\sum_{k=1}^{K}\alpha_k,
\]

the backend must calculate

\[
K(\mathbf{x};\boldsymbol{\alpha})
=
\log\Gamma(\alpha_0)
-\log\Gamma(\alpha_0+N)
+\sum_{k=1}^{K}
\left[
\log\Gamma(\alpha_k+x_k)-\log\Gamma(\alpha_k)
\right].
\]

For a count matrix

\[
X\in\mathbb{N}_0^{R\times K},
\]

each row is an independent observation window. The backend result is

\[
K(X;\boldsymbol{\alpha})
=
\sum_{r=1}^{R}K(\mathbf{x}_r;\boldsymbol{\alpha}).
\]

The rows must not be added together before evaluation. Aggregating them would
describe one different observation and would not match the existing backend
contract.

The backend must not include the multinomial coefficient

\[
\log\frac{N!}{\prod_{k=1}^{K}x_k!}.
\]

That count-only term is already added outside the backend when a complete
window probability is needed. Keeping it outside preserves identical semantics
for LM and SciPy and keeps model-fitting convergence values separate from
complete anomaly scores.

\section{Conversion to the Native Parameterization}

The Python application works with \(\boldsymbol{\alpha}\). The existing C
implementation works with category probabilities and the dispersion parameter
\(\psi\). The adapter will perform the conversion once per matrix call:

\[
\alpha_0=\sum_{k=1}^{K}\alpha_k,
\qquad
p_k=\frac{\alpha_k}{\alpha_0},
\qquad
\psi=\frac{1}{\alpha_0}.
\]

Inside the existing C code,

\[
\frac{\psi}{p_k}=\frac{1}{\alpha_k},
\]

which is the reciprocal parameter expected by the LM paired log-gamma
difference calculation.

This conversion belongs in the Python native adapter. It should not be
duplicated in training, calibration, or scoring code.

\section{Implementation Decisions}

The following decisions keep the integration small and consistent with the
current application.

\begin{enumerate}

\item Modify only the runtime project copy in
\path{src/lm_idnet/algorithms/native/}. Keep \path{lm-algorithm/} unchanged as
the reference capsule.

\item Add the minimum application-facing functions directly to
\path{LM_time_lib.c} and declare them in \path{LM_hor_shift.h}. A separate C
framework or a second copy of the numerical implementation is unnecessary.

\item Send an entire row-major count matrix through one Python-to-C call. The
C wrapper will evaluate and sum the rows independently.

\item Retain the current native global state and serialize access with one
Python lock. A native context structure is unnecessary for the current
single-stream command-line pipeline.

\item Initialize the two coefficient tables needed by the log-likelihood
calculation using absolute paths supplied by Python. The Python process will
not change its working directory.

\item Reuse the existing Python input validation before entering C. The native
wrapper will still reject invalid dimensions, failed allocations, missing
coefficient state, and non-finite results defensively.

\item Keep SciPy as the numerical reference backend. Do not remove it and do
not silently fall back to it when LM fails.

\item Keep the current artifact schema unchanged in this task. The backend name
is already recorded in the model, while precision comes from the validated
configuration used for the run. Persisting precision inside standalone model
artifacts would require a separate artifact-versioning decision.

\end{enumerate}

\section{Step-by-Step Work Plan}

\subsection{Step 1: Record a clean baseline}

Before editing code, run the existing focused tests:

\begin{verbatim}
python -m pytest tests/test_likelihood.py tests/test_estimator_factory.py
\end{verbatim}

Then run the normal suite:

\begin{verbatim}
python -m pytest
\end{verbatim}

The purpose is to separate existing failures from failures introduced by the
integration. The current placeholder test is expected to pass at this stage
because it verifies that LM is unavailable.

\textbf{Checkpoint:} record the commands used and confirm that the worktree has
no unrelated changes that could be overwritten.

\subsection{Step 2: Establish the native source boundary}

Use the files under \path{src/lm_idnet/algorithms/native/} as the runtime source
of truth. Compare them with the corresponding files under
\path{lm-algorithm/code/} before editing so that project-specific changes can be
identified later.

Only the log-likelihood portion of the native code is required. Do not import
the experimental Python scripts, dataset files, estimator implementation,
digamma routines, trigamma routines, or benchmark driver from the capsule.
Minka's estimator already exists in the application.

\textbf{Checkpoint:} identify the existing functions used by the likelihood
path: coefficient initialization, parameter initialization,
\texttt{LM\_logL}, error-bound selection, and \texttt{loggamma\_LM}.

\subsection{Step 3: Make zero counts mathematically correct}

Add an early case to the paired log-gamma difference function:

\begin{verbatim}
if (y == 0) {
    return 0.0;
}
\end{verbatim}

This implements

\[
\log\Gamma(\alpha+0)-\log\Gamma(\alpha)=0.
\]

It is required because protocol categories can be absent from a window and a
valid silent window can contain only zeros. Existing special cases for counts
one and two remain unchanged.

\textbf{Checkpoint:} a zero category contributes zero, and an all-zero row
produces a row kernel of zero.

\subsection{Step 4: Load coefficient tables by explicit path}

Replace the current-working-directory assumption with an initializer similar
to:

\begin{verbatim}
int initBern_logL_paths(
    const char *bernoulli_path,
    const char *error_path
);
\end{verbatim}

Only two supplied files are needed for the log-likelihood path:

\begin{itemize}

\item \path{bernreal-norm-100.txt}; and

\item \path{err_coeff-100.txt}.

\end{itemize}

The initializer must open both files, allocate both 100-element tables, verify
all 100 reads from each file, and return a nonzero status if any operation
fails. New tables should replace existing tables only after both files have
been read successfully. Previously allocated tables must be released before
replacement.

Python will build the paths from the location of \path{lm_native.py}; it will
not assume that the command was launched from the repository root.

\textbf{Checkpoint:} initialization succeeds after changing the process to an
unrelated working directory, and a missing or malformed table produces a clear
Python exception.

\subsection{Step 5: Define and clean up native workspace ownership}

The existing native parameter initializer allocates arrays for counts,
probabilities, reciprocal parameters, and logarithms. Add one cleanup helper
that frees these arrays and sets their pointers to \texttt{NULL}.

Call the cleanup helper:

\begin{itemize}

\item before replacing an existing workspace;

\item after the last row of a matrix has been evaluated; and

\item before returning from a calculation failure.

\end{itemize}

Check every allocation before writing through its pointer. A failed allocation
must be reported to the matrix wrapper; it must not become a segmentation
fault.

Coefficient tables have a longer lifetime than per-row workspace. They may
remain initialized for the lifetime of the loaded backend, but repeated backend
construction must not leak the previous tables.

\textbf{Checkpoint:} repeated calls with different row and category counts do
not reuse stale data and do not continually increase native allocations.

\subsection{Step 6: Add one matrix-level C entry point}

Expose a function that receives the whole contiguous count matrix, its shape,
the converted probability vector, \(\psi\), and the requested precision. A
minimal interface may return the result directly and use \texttt{NAN} for a
native failure:

\begin{verbatim}
double loggamma_LM_matrix(
    int precision_digits,
    const double *counts,
    long rows,
    int categories,
    const double *probabilities,
    double psi
);
\end{verbatim}

The function must perform the following work in order:

\begin{enumerate}

\item reject nonpositive row counts, fewer than two categories, null pointers,
or uninitialized coefficient tables;

\item initialize the native parameters from the first row;

\item evaluate that row with the existing LM numerical routine;

\item repeat initialization and evaluation for each remaining row;

\item add the row results to one \texttt{double};

\item stop on a precision, allocation, or non-finite-result failure;

\item release per-row workspace; and

\item return the finite sum on success or \texttt{NAN} on failure.

\end{enumerate}

Any existing native path that prints and exits must not be reachable through
this library entry point. Existing special sentinel values must be converted
to a failure before returning to Python.

This step changes only the application boundary and lifecycle around the LM
calculation. The expansion, error coefficients, error bound, term selection,
and horizontal-shift mathematics remain unchanged.

\textbf{Checkpoint:} one C call can evaluate a multirow matrix without
aggregating its rows.

\subsection{Step 7: Rebuild the shared library}

Compile the edited source locally from the native directory:

\begin{verbatim}
cd src/lm_idnet/algorithms/native
gcc -O3 -fPIC -shared LM_time_lib.c -lm -o LM_time_lib.so
\end{verbatim}

The existing Docker entry point already performs the equivalent compilation
through a temporary output file before each command. Verify that the new
symbols are present by loading the library through a small Python check or by
running the focused likelihood tests.

Do not add a new build framework or dependency for this integration. The
current Linux and Docker execution path already has the required compiler and
linker configuration.

\textbf{Checkpoint:} the shared library builds without warnings that indicate
an invalid declaration, missing symbol, or missing math-library linkage.

\subsection{Step 8: Implement the Python native adapter}

Implement \path{src/lm_idnet/algorithms/lm_native.py} as the only module that
knows about \texttt{ctypes}.

The module should contain:

\begin{itemize}

\item a native-directory path derived from \texttt{\_\_file\_\_};

\item the common \texttt{double *} type declaration;

\item one process-wide \texttt{threading.Lock}; and

\item a small \texttt{NativeLmKernel} class.

\end{itemize}

During construction, \texttt{NativeLmKernel} will:

\begin{enumerate}

\item validate and retain the requested precision digits;

\item load \path{LM_time_lib.so} from the native directory;

\item declare \texttt{argtypes} and \texttt{restype} for the two public C
functions;

\item pass absolute paths for the two coefficient files while holding the
global lock; and

\item raise \texttt{CommandUnavailableError} if the library or tables cannot
be loaded.

\end{enumerate}

For each calculation, it will:

\begin{enumerate}

\item make the count matrix contiguous with NumPy \texttt{float64} storage;

\item calculate concentration, probabilities, and \(\psi\) from
\(\boldsymbol{\alpha}\);

\item make the probability vector contiguous;

\item call the matrix-level native function while holding the global lock;

\item convert the result to a Python \texttt{float}; and

\item raise \texttt{NumericalPrecisionError} if the result is not finite.

\end{enumerate}

The lock must cover native initialization and every native calculation because
the loaded library stores dimensions, parameters, coefficient pointers, and
workspace in global variables.

\textbf{Checkpoint:} no other Python module imports \texttt{ctypes} or calls a
native LM function directly.

\subsection{Step 9: Replace the Python placeholder}

Update \path{src/lm_idnet/algorithms/log_likelihood.py} so that
\texttt{LmLogLikelihood} owns a \texttt{NativeLmKernel}.

Its \texttt{calculate} method should:

\begin{enumerate}

\item convert counts and alpha to NumPy arrays;

\item reuse the same validation currently used by
\texttt{ScipyLogLikelihood}; and

\item delegate the valid arrays to the native adapter.

\end{enumerate}

Do not duplicate the dimension, integer-count, non-negativity, finiteness, or
positive-alpha checks. Move the validator to the common base only if doing so
makes the code shorter; otherwise reuse the existing static validator.

Allow \texttt{LmLogLikelihood} and \texttt{initialize\_log\_likelihood} to
receive \texttt{precision\_digits}, with the current configured default of six
digits. The SciPy constructor does not need this setting.

\textbf{Checkpoint:} \texttt{initialize\_log\_likelihood("lm", 6)} returns a
working \texttt{LmLogLikelihood}, while an unknown backend still raises the
existing validation error.

\subsection{Step 10: Pass precision through every construction path}

Update \path{src/lm_idnet/algorithms/estimator_factory.py} so that
\texttt{create\_estimator} accepts the precision setting and forwards it to
\texttt{initialize\_log\_likelihood}.

Then pass \texttt{config.precision\_digits} at every current application call
site:

\begin{itemize}

\item model fitting in \path{src/lm_idnet/models/modeling_stage.py};

\item daily diagnostics in \path{src/lm_idnet/cli.py};

\item threshold calibration in
\path{src/lm_idnet/models/calibration_stage.py}; and

\item window scoring in \path{src/lm_idnet/models/scoring_stage.py}.

\end{itemize}

Calibration and scoring must continue to take the backend name from the model
artifact, not from the current estimator selection. They use the current
validated configuration only for the precision value.

Keep a default precision argument where it preserves existing direct calls and
tests. Do not create a second precision setting inside the estimator section.

\textbf{Checkpoint:} selecting LM through configuration reaches the same
backend in diagnostics, fitting, calibration, and scoring.

\subsection{Step 11: Replace the placeholder test with numerical tests}

Update \path{tests/test_likelihood.py}. Remove the test that expects the LM
backend to be unavailable and add the following checks.

\begin{enumerate}

\item \textbf{Known small matrix.} Compare LM and SciPy on a small manually
understandable count matrix.

\item \textbf{Zero categories.} Include rows in which individual categories
are zero.

\item \textbf{Silent row.} Include an all-zero row and confirm that it changes
the matrix kernel by exactly zero.

\item \textbf{Independent rows.} Compare a multirow call with the sum of calls
made one row at a time. This prevents accidental row aggregation.

\item \textbf{Seeded stress sample.} Generate deterministic matrices with
different row counts, counts up to approximately \(10^5\), and alpha values
spanning small and large magnitudes. Compare LM with SciPy using an absolute
tolerance of \(10^{-6}\) for six requested digits.

\item \textbf{Repeated dimensions.} Reuse one backend for successive matrices
of different shapes to expose stale native state or cleanup errors.

\item \textbf{Working-directory independence.} Change to a temporary directory
before constructing the backend and confirm that coefficient initialization
still succeeds.

\item \textbf{Existing validation.} Run invalid dimensions, negative or
fractional counts, and invalid alpha values through both backends and confirm
that both reject them consistently.

\end{enumerate}

Keep the random seed fixed and report the largest observed absolute difference
if a test fails. Do not use randomly changing test data.

Update \path{tests/test_estimator_factory.py} only as needed to verify that the
selected backend and precision reach the estimator. Update downstream tests
whose expected backend changes only when the corresponding test configuration
is deliberately changed to LM.

\textbf{Checkpoint:} the numerical tests would fail if rows were aggregated,
zero counts were mishandled, coefficient paths were relative, or native state
leaked between calls.

\subsection{Step 12: Run the focused and complete test suites}

After rebuilding the shared library, run:

\begin{verbatim}
python -m pytest tests/test_likelihood.py tests/test_estimator_factory.py
python -m pytest -m numerical
python -m pytest
\end{verbatim}

If a test fails, determine whether the failure is caused by the C calculation,
the parameter conversion, the row-wise wrapper, precision propagation, or the
test expectation. Do not weaken the tolerance merely to make an unexplained
difference pass.

Run the same focused tests inside Docker so that the compiled environment used
by normal commands is also covered.

\textbf{Checkpoint:} all existing SciPy behavior remains valid and every new LM
test passes locally and in Docker.

\subsection{Step 13: Exercise the complete pipeline}

Use a copied configuration or temporary artifact paths for the first LM run so
that accepted SciPy artifacts are not overwritten during validation. Set:

\begin{verbatim}
"log_likelihood_backend": "lm"
\end{verbatim}

Then run, in order:

\begin{verbatim}
lm-idnet diagnose  --config <lm-config>
lm-idnet train     --config <lm-config>
lm-idnet calibrate --config <lm-config>
lm-idnet score     --config <lm-config>
\end{verbatim}

Confirm that:

\begin{itemize}

\item model fitting converges;

\item initial and final kernel values are finite;

\item the model records \texttt{lm} as its backend;

\item calibration uses the LM backend recorded by the model;

\item scoring uses that same backend;

\item silent windows remain valid; and

\item all saved artifacts pass their existing schema validation.

\end{itemize}

Run the same data once with SciPy and compare fitted alpha values, thresholds,
and anomaly decisions. Small numerical differences are acceptable within the
requested precision. Unexplained convergence or decision changes are not.

\textbf{Checkpoint:} the complete LM path works without a SciPy fallback and
without changing the mathematical meaning of saved scores.

\subsection{Step 14: Measure performance without assuming a speedup}

Record at least:

\begin{itemize}

\item total model-fitting time;

\item likelihood time for a representative matrix;

\item single-window scoring latency; and

\item the equivalent SciPy measurements on the same data and machine.

\end{itemize}

The acceptance criterion is that runtime remains practical for the pipeline.
The integration must not claim that LM is faster than vectorized SciPy unless
the measurements demonstrate it. The primary reason for LM is controlled
evaluation of paired log-gamma differences in difficult numerical regimes.

\textbf{Checkpoint:} numerical and detection comparisons are recorded together
with runtime measurements.

\subsection{Step 15: Enable LM in the selected research configurations}

Only after the numerical and end-to-end checks pass, change the intended
versioned experiment configurations from \texttt{scipy} to \texttt{lm}. Keep
SciPy available as the reference implementation and keep configuration
validation restricted to the two known names.

Regenerate affected model, threshold, and event artifacts through their normal
pipeline commands. Do not edit numerical artifact values by hand. Review the
diff before accepting regenerated artifacts.

\textbf{Checkpoint:} each regenerated model records the LM backend and every
dependent threshold was calibrated from that model.

\subsection{Step 16: Update user-facing documentation}

Update the following documentation after the implementation is proven:

\begin{itemize}

\item \path{docs/likelihood-semantics.md}: describe LM as the working native
implementation and SciPy as the reference;

\item \path{docs/configuration.md}: remove the placeholder warning and explain
the two backend choices;

\item \path{README.md}: change it only if the actual build or run instructions
differ from the existing Docker compilation flow; and

\item the experiment report: record numerical agreement, detection agreement,
runtime, compiler environment, and configured precision.

\end{itemize}

Preserve the existing native licence and author attribution. Document that the
project copy contains integration and lifecycle changes, while making no claim
that those changes alter the underlying LM method.

\textbf{Checkpoint:} no documentation still describes \texttt{lm} as a
placeholder.

\section{Files Expected to Change}

The main implementation should be limited to the following files:

\begin{itemize}

\item \path{src/lm_idnet/algorithms/native/LM_time_lib.c};

\item \path{src/lm_idnet/algorithms/native/LM_hor_shift.h};

\item \path{src/lm_idnet/algorithms/lm_native.py};

\item \path{src/lm_idnet/algorithms/log_likelihood.py};

\item \path{src/lm_idnet/algorithms/estimator_factory.py};

\item \path{src/lm_idnet/models/modeling_stage.py};

\item \path{src/lm_idnet/models/calibration_stage.py};

\item \path{src/lm_idnet/models/scoring_stage.py};

\item \path{src/lm_idnet/cli.py};

\item \path{tests/test_likelihood.py};

\item \path{tests/test_estimator_factory.py}, if precision wiring needs an
explicit assertion;

\item \path{docs/likelihood-semantics.md}; and

\item \path{docs/configuration.md}.

\end{itemize}

The Docker entry point, artifact schemas, estimator mathematics, and
Dirichlet--Multinomial score composition should not require changes. Dataset
configuration and generated artifact files change only after validation when
LM is selected for an experiment.

\section{Important Failure Modes to Guard Against}

\begin{itemize}

\item \textbf{Aggregating rows:} produces the likelihood of one combined count
vector instead of the sum of independent-window kernels.

\item \textbf{Including the multinomial coefficient twice:} changes anomaly
scores because the application already adds this term outside the backend.

\item \textbf{Incorrect parameter conversion:} passing \(\alpha_0\) instead of
\(1/\alpha_0\), or passing alpha directly as probabilities, changes the model.

\item \textbf{Relative coefficient paths:} makes commands work only from one
directory.

\item \textbf{Missing zero-count handling:} breaks common sparse windows and
silent observations.

\item \textbf{Unlocked global state:} permits one call to overwrite another
call's dimensions or buffers.

\item \textbf{Stale shared library:} Python may load an old binary after the C
source changes. Rebuild before every native test.

\item \textbf{Silent fallback:} hiding an LM error by running SciPy would make
the recorded backend provenance false.

\item \textbf{Sentinel values treated as likelihoods:} native error markers
must be converted to explicit failures before the result reaches the pipeline.

\end{itemize}

\section{Work Deliberately Outside This Task}

The following work should not be added unless a demonstrated requirement
appears:

\begin{itemize}

\item rewriting the LM calculation in Python;

\item replacing Minka's fixed-point estimator;

\item integrating unused digamma or trigamma native routines;

\item introducing a native context-object API for concurrent calls;

\item adding a new foreign-function or build dependency;

\item designing cross-platform binary distribution outside the current Linux
and Docker target;

\item removing the SciPy backend;

\item automatically falling back from LM to SciPy;

\item changing the model-artifact schema solely to store precision; or

\item modifying the LM expansion, coefficients, error analysis, or horizontal
shift algorithm.

\end{itemize}

These items can be reconsidered only if testing, deployment, or measured
concurrency requirements show that the smaller design is insufficient.

\section{Final Acceptance Checklist}

Before declaring the backend complete, verify each item below:

\begin{enumerate}

\item The native library compiles successfully in the supported Docker image.

\item LM and SciPy implement the same kernel semantics.

\item Multirow results equal the sum of individual-row results.

\item Zero categories and silent rows return correct finite values.

\item Coefficient lookup works from an arbitrary working directory.

\item Repeated calls with different shapes remain correct.

\item Requested precision reaches the C function from every pipeline command.

\item Native loading, coefficient, allocation, precision, and numerical errors
do not terminate the Python process.

\item Seeded numerical comparisons satisfy the stated tolerance.

\item The focused, numerical, and complete test suites pass.

\item Training, calibration, and scoring complete with LM and produce valid
artifacts.

\item LM and SciPy experiment results have been compared and any material
difference has been explained.

\item Documentation no longer calls LM a placeholder and does not claim an
unmeasured performance advantage.

\item The original licence and attribution remain with the native source.

\end{enumerate}

\end{document}
